% Milestone 2 benchmark results — paste into a larger report or compile standalone.
% Build: pdflatex milestone2_results.tex
\documentclass[11pt]{article}

\usepackage[margin=1in]{geometry}
\usepackage{booktabs}
\usepackage{multirow}
\usepackage{graphicx}
\usepackage[hidelinks]{hyperref}
\usepackage{caption}

\captionsetup{font=small, labelfont=bf}

\title{Milestone 2}
\author{Sarosh Khan and Naomi Chien}
\date{5/8/2026}

\begin{document}
\maketitle

\noindent\textbf{Environment.}
Model: \texttt{Qwen/Qwen3-8B};\\
GPU: NVIDIA L4;\\
Paged runs reserve KV memory with \texttt{--mem-fraction-static 0.85} unless noted.

% ---------------------------------------------------------------------------
\section{Design and implementation}

Milestone~1 decoded with \texttt{torch.cat}-grown KV caches, which reallocates, fragments GPU memory, and breaks pointer stability under load.
We replace that with a fixed-budget pool (\texttt{KVMemoryPool}): page-major \(K,V\) tensors per layer, a free-list of physical pages, per-request page tables (\texttt{PagedRequestState}), and admission via peak token count (\texttt{paged\_kv\_can\_admit}); \texttt{from\_budget} sizes the pool from a VRAM fraction.
Prefill is packed varlen (\texttt{flash\_attn\_varlen\_func}); decode indexes the pool via \texttt{flash\_attn\_with\_kvcache} and \texttt{block\_table} when \texttt{page\_size} is divisible by~256---otherwise pooled KV may be gathered in \texttt{\_build\_batched\_kv\_from\_pool}.
\texttt{Attention.forward} selects among these branches (else SDPA plus the legacy cache path).

For launch overhead we \texttt{torch.compile} each layer's SwiGLU \texttt{mlp} (\texttt{mode=reduce-overhead}, CLI \texttt{--torch-compile}) with synthetic batch warmup (and CUDA-graph step markers when available).

% ---------------------------------------------------------------------------
\section{Experiments}

\subsection{Accuracy (MMLU)}

\noindent Commands:
\begin{verbatim}
python -m miniengine --model Qwen/Qwen3-8B --mode batched

python -m miniengine --model Qwen/Qwen3-8B --mode paged \
  --mem-fraction-static 0.85 --page-size 256

python -m miniengine --model Qwen/Qwen3-8B --mode paged \
  --mem-fraction-static 0.85 --page-size 256 \
  --torch-compile --torch-compile-target mlp
\end{verbatim}

\begin{table}[htbp]
  \centering
  \caption{MMLU accuracy (100 samples).}
  \label{tab:mmlu}
  \begin{tabular}{@{}lccc@{}}
    \toprule
    Configuration & Accuracy & Correct / Samples & Avg latency \\
    \midrule
    Batched (Milestone~1 baseline) & 59.0\% & 59 / 100 & 2.07\,s \\
    Paged (\texttt{page-size}=256) & 60.0\% & 60 / 100 & 1.74\,s \\
    Paged + \texttt{torch.compile} (MLP) & 60.0\% & 60 / 100 & 1.55\,s \\
    \bottomrule
  \end{tabular}
\end{table}

% ---------------------------------------------------------------------------
\subsection{Throughput: batched baseline vs.\ paged}

\noindent We sweep concurrency at fixed prompt/generation lengths. We report time-to-first-token (TTFT), end-to-end completion time, time per output token (TPOT), and aggregate generation throughput (GenTok/s).

\noindent Benchmark:
\begin{verbatim}
python3 -m benchmark.bench_serving --input-len 512 --output-len 256 \
  --concurrencies 1,4,8,16
\end{verbatim}

\noindent Engine (batched):
\begin{verbatim}
python -m miniengine --model Qwen/Qwen3-8B --mode batched
\end{verbatim}

\noindent Engine (paged):
\begin{verbatim}
python -m miniengine --model Qwen/Qwen3-8B --mode paged \
  --mem-fraction-static 0.85 --page-size 256
\end{verbatim}

\begin{table}[htbp]
  \centering
  \caption{Batched vs.\ paged \texttt{bench\_serving} (512 input / 256 output tokens). TTFT and completion in ms; TPOT in ms/token.}
  \label{tab:batched-paged}
  \resizebox{\textwidth}{!}{%
  \begin{tabular}{@{}rcccccccc@{}}
    \toprule
    Mode & Conc & TTFT$_{p50}$ & TTFT$_{p99}$ & Compl$_{p50}$ & Compl$_{p99}$ & TPOT$_{p50}$ & TPOT$_{p99}$ & GenTok/s \\
    \midrule
    \multirow{5}{*}{Batched}
      &  1 & 220 & 542 & 14019 & 26129 & 70.8 & 111.3 & 13 \\
      &  4 & 312 & 792 & 18126 & 22441 & 93.6 & 95.7 & 39 \\
      &  8 & 336 & 1651 & 22274 & 29445 & 121.0 & 125.1 & 53 \\
      & 16 & 396 & 3295 & 31251 & 45163 & 181.1 & 396.9 & 68 \\
      & 32 & 20985 & 41631 & 59891 & 103175 & 231.0 & 425.4 & 72 \\
    \midrule
    \multirow{5}{*}{Paged}
      &  1 & 221 & 976 & 13291 & 15817 & 66.8 & 67.0 & 15 \\
      &  4 & 369 & 796 & 14726 & 17465 & 74.3 & 75.9 & 48 \\
      &  8 & 384 & 1297 & 14845 & 19916 & 79.4 & 85.5 & 79 \\
      & 16 & 311 & 2571 & 16516 & 23591 & 88.0 & 113.3 & 129 \\
      & 32 & 10481 & 22769 & 29341 & 50076 & 117.0 & 294.0 & 151 \\
    \bottomrule
  \end{tabular}%
  }
\end{table}

% ---------------------------------------------------------------------------
\subsection{Throughput: paged vs.\ paged + \texttt{torch.compile}}

\noindent Same serving benchmark, but with randomized prompts and fixed request count to better expose steady-state decode overhead. We compare paged eager vs.\ paged with compiled MLP submodules.

\noindent Benchmark:
\begin{verbatim}
python3 -m benchmark.bench_serving --input-len 512 --output-len 256 \
  --randomness 1.0 --concurrencies 8,16,32 --num-requests 64
\end{verbatim}

\noindent Engine (paged baseline):
\begin{verbatim}
python -m miniengine --model Qwen/Qwen3-8B --mode paged \
  --mem-fraction-static 0.85 --page-size 256
\end{verbatim}

\noindent Engine (paged + compile):
\begin{verbatim}
python -m miniengine --model Qwen/Qwen3-8B --mode paged \
  --mem-fraction-static 0.85 --page-size 256 \
  --torch-compile --torch-compile-target mlp
\end{verbatim}

\begin{table}[htbp]
  \centering
  \caption{Effect of \texttt{torch.compile} (MLP target) on paged serving.}
  \label{tab:compile}
  \resizebox{\textwidth}{!}{%
  \begin{tabular}{@{}rcccccccc@{}}
    \toprule
    Configuration & Conc & TTFT$_{p50}$ & TTFT$_{p99}$ & Compl$_{p50}$ & Compl$_{p99}$ & TPOT$_{p50}$ & TPOT$_{p99}$ & GenTok/s \\
    \midrule
    \multirow{3}{*}{Paged}
      &  8 & 346 & 3620 & 21762 & 23187 & 83.6 & 105.8 & 90 \\
      & 16 & 2910 & 16273 & 24207 & 38861 & 89.8 & 100.4 & 136 \\
      & 32 & 15699 & 35190 & 42104 & 77424 & 123.1 & 442.1 & 134 \\
    \midrule
    \multirow{3}{*}{\shortstack{Paged +\\ \texttt{torch.compile}}}
      &  8 & 343 & 1646 & 21572 & 22298 & 83.4 & 105.4 & 91 \\
      & 16 & 1908 & 3202 & 24891 & 25653 & 90.1 & 113.0 & 150 \\
      & 32 & 12573 & 26922 & 37645 & 57951 & 119.5 & 429.7 & 152 \\
    \bottomrule
  \end{tabular}%
  }
\end{table}

% ---------------------------------------------------------------------------
\subsection{Page-size comparison}

\noindent Same workload as the compile comparison, varying only the KV page size. Smaller pages reduce tail fragmentation; larger pages reduce page-table indirection.

\noindent Benchmark (same as previous section):
\begin{verbatim}
python3 -m benchmark.bench_serving --input-len 512 --output-len 256 \
  --randomness 1.0 --concurrencies 8,16,32 --num-requests 64
\end{verbatim}

\noindent Engines compared:
\begin{verbatim}
python -m miniengine --model Qwen/Qwen3-8B --mode paged \
  --mem-fraction-static 0.85 --page-size 256

python -m miniengine --model Qwen/Qwen3-8B --mode paged \
  --mem-fraction-static 0.85 --page-size 1024
\end{verbatim}

\begin{table}[htbp]
  \centering
  \caption{\texttt{bench\_serving} for two KV page sizes.}
  \label{tab:pagesize}
  \resizebox{\textwidth}{!}{%
  \begin{tabular}{@{}rcccccccc@{}}
    \toprule
    Page size & Conc & TTFT$_{p50}$ & TTFT$_{p99}$ & Compl$_{p50}$ & Compl$_{p99}$ & TPOT$_{p50}$ & TPOT$_{p99}$ & GenTok/s \\
    \midrule
    \multirow{3}{*}{256}
      &  8 & 346 & 3620 & 21762 & 23187 & 83.6 & 105.8 & 90 \\
      & 16 & 2910 & 16273 & 24207 & 38861 & 89.8 & 100.4 & 136 \\
      & 32 & 15699 & 35190 & 42104 & 77424 & 123.1 & 442.1 & 134 \\
    \midrule
    \multirow{3}{*}{1024}
      &  8 & 286 & 3179 & 15806 & 21383 & 82.9 & 86.5 & 90 \\
      & 16 & 2711 & 11665 & 19794 & 32116 & 90.2 & 98.6 & 131 \\
      & 32 & 14171 & 28943 & 33296 & 59648 & 119.9 & 304.6 & 130 \\
    \bottomrule
  \end{tabular}%
  }
\end{table}

% ---------------------------------------------------------------------------
\section{Sources of performance benefits and analysis}

\paragraph{Batched baseline vs.\ paged (Table~\ref{tab:batched-paged}).}
The primary gain comes from avoiding dynamic KV growth via \texttt{torch.cat} during decode. With a paged pool, KV appends are implemented as indexed writes into pre-allocated pages, which reduces reallocations/fragmentation and keeps memory stable under high concurrency. Empirically this yields higher throughput (conc.\ 16: 68 \(\rightarrow\) 129 GenTok/s; conc.\ 32: 72 \(\rightarrow\) 151) and improved tail behavior (conc.\ 16 TPOT${}_{p99}$ \(\sim\)397 \(\rightarrow\) \(\sim\)113 ms/token). At low concurrency, TTFT differences are small and more sensitive to warmup and scheduling overheads.

\paragraph{\texttt{torch.compile} on the MLP (Table~\ref{tab:compile}).}
We compile only the per-layer MLP, since it is a stable region executed at every decode step. This primarily reduces kernel-launch overhead and provides a modest but consistent throughput increase at higher concurrency (conc.\ 16: 136 \(\rightarrow\) 150 GenTok/s; conc.\ 32: 134 \(\rightarrow\) 152). TPOT changes are small because attention remains a major component of per-token cost.

\paragraph{Accuracy sweep (Table~\ref{tab:mmlu}).}
On 100 MMLU samples, 59 vs.\ 60 correct is within expected sampling noise, so we do not interpret the difference as meaningful. The key result is that paging and \texttt{torch.compile} improve performance without an observable accuracy regression on this small check. Average latency per example improves (2.07 s \(\rightarrow\) 1.74 s \(\rightarrow\) 1.55 s), consistent with cheaper KV maintenance and reduced launch overhead.

\paragraph{Page size 256 vs.\ 1024 (Table~\ref{tab:pagesize}).}
GenTok/s is similar at conc.\ 8, while smaller pages are slightly faster at conc.\ 16/32; larger pages sometimes reduce median completion time. This matches the expected tradeoff: smaller pages reduce tail fragmentation at the cost of more bookkeeping, while larger pages reduce metadata/indirection overhead but may waste more KV at sequence tails.

\end{document}
